\documentclass[journal]{IEEEtran}
% IEEE TNSM / Computers & Security submission draft
% Compile: pdflatex main.tex (or upload the paper/ folder to Overleaf)

\usepackage{amsmath,amssymb,amsthm}
\usepackage{booktabs}
\usepackage{graphicx}
\usepackage{url}
\usepackage[hidelinks]{hyperref}

\newtheorem{theorem}{Theorem}
\newtheorem{proposition}{Proposition}
\newtheorem{corollary}{Corollary}
\newtheorem{definition}{Definition}

\newcommand{\TPR}{\mathrm{TPR}}
\newcommand{\FPR}{\mathrm{FPR}}
\newcommand{\KS}{\mathrm{KS}}
\newcommand{\OR}{\mathrm{OR}}

\begin{document}

\title{Anomalous Is Not Malicious: A Provable Performance Ceiling for
Unsupervised NIDS and a Two-Stage Architecture for Zero-Day Intrusion Detection}

\author{Anurag~Sharma%
\thanks{A.~Sharma is with the Advanced Data Processing Research Institute
(ADRIN), Indian Space Research Organisation (ISRO), Hyderabad, India
(e-mail: anuragadrin@gmail.com).}}

\markboth{Working draft --- not yet submitted}%
{Sharma: Anomalous Is Not Malicious}

\maketitle

\begin{abstract}
Unsupervised network intrusion detection systems (NIDS) consistently plateau at
55--70\% recall on real-world benchmarks, a degradation widely attributed to
parameter misconfiguration. We prove this ceiling is \emph{structural}: for any
threshold-based detector, the maximum achievable Youden index
($\TPR-\FPR$) equals the Kolmogorov--Smirnov (KS) distance between the
class-conditional anomaly-score distributions (Theorem~\ref{thm:ks}). On the
CTU-13 botnet dataset, Isolation Forest attains $\KS=0.445$; no threshold
calibration can exceed this bound. The root cause is the
\emph{anomaly--malice mismatch}: unsupervised detectors measure statistical
rarity, whereas stealthy attacks are engineered to appear statistically normal.
Replacing the anomaly notion --- isolation by reconstruction, via an
autoencoder trained exclusively on normal traffic --- raises the ceiling to
$\KS=0.860$ ($F_1$: $0.617\!\to\!0.928$), and the deployed detector operates
\emph{at} its theoretical ceiling (balanced accuracy $0.929$ vs.\ bound
$0.930$). Cross-dataset evaluation on UNSW-NB15, however, shows no single
anomaly notion generalises across attack families. We therefore derive a
closed-form law for OR-fusion of conditionally independent detectors
(Theorem~\ref{thm:fusion}): miss rates multiply. The law predicts the fused
recall of our dual-head detector (flow autoencoder $\cup$ temporal autoencoder)
to three decimal places (predicted $0.488$, measured $0.491$;
miss-redundancy $\rho_{\mathrm{miss}}=0.994$), lifting Stage-1 recall from
$8.6\%$ to $49.1\%$ and, on hidden attack families, Backdoor from $8\%$ to
$41\%$ and Analysis from $3\%$ to $35\%$. Finally, we replace the arbitrary
percentile threshold with an extreme-value-theoretic threshold
(Theorem~\ref{thm:evt}): a generalized Pareto tail fit
($\hat\xi = 1.12$, infinite variance) yields calibrated false-positive rates
where percentile thresholds overshoot by $3.8\times$. The resulting two-stage
architecture --- fused zero-day net, gradient-boosted known-attack classifier,
and a zero-day candidate queue --- achieves $F_1=0.945$ on CTU-13 with 87\%
false-positive reduction. In a controlled simulation withholding four attack
families from the classifier entirely, the dual-head pipeline raises end-to-end
zero-day routing of Backdoor from 5.0\% to 27.8\% and Analysis from 2.2\% to
23.8\%, with zero-day queue precision of 86.7\%; of all hidden-family flows
flagged by Stage~1, fewer than 1\% are cleared as benign --- the remainder
reach an analyst either through the queue or as (family-misattributed) attack
alerts.
\end{abstract}

\begin{IEEEkeywords}
Network intrusion detection, anomaly detection, zero-day attacks,
Kolmogorov--Smirnov distance, extreme value theory, detector fusion,
Isolation Forest, autoencoder, CTU-13, UNSW-NB15.
\end{IEEEkeywords}

\section{Introduction}
\IEEEPARstart{S}{ignature-based} intrusion detection fails categorically
against zero-day attacks: no signature exists for an attack no one has seen.
Anomaly-based NIDS promise the alternative --- model normal traffic, flag
deviations --- yet two decades of published results place unsupervised recall
stubbornly in the 55--70\% band \cite{garcia2014,moustafa2016}. The standard
prescriptions (tune the contamination parameter, add features, deepen the
model) do not move this number. This paper explains why, proves the limitation
formally, and presents a validated architecture that circumvents it.

Our contributions:
\begin{enumerate}
  \item \textbf{A provable ceiling (Sec.~\ref{sec:theory}).} For threshold
  detectors, $\sup_t J(t) = \KS$ exactly. Corollaries give a balanced-accuracy
  ceiling $(1+\KS)/2$ and a recall bound at any false-alarm budget. On CTU-13
  our autoencoder \emph{attains} its ceiling, proving threshold optimisation is
  exhausted and only representation change can help.
  \item \textbf{A fusion law (Sec.~\ref{sec:theory}).} Under conditional
  independence, OR-fused detector misses multiply. The law is verified on real
  traffic to three decimal places and supplies a closed-form criterion for when
  fusing two detectors is beneficial.
  \item \textbf{Parameter-free thresholds (Sec.~\ref{sec:theory}).} A
  peaks-over-threshold GPD fit converts any target false-positive budget into a
  calibrated threshold, including budgets below the empirical resolution of the
  training data.
  \item \textbf{A validated two-stage system (Sec.~\ref{sec:system}).}
  Dual-head Stage~1 (flow + temporal autoencoders, fused per the law), a
  supervised Stage~2 over the flagged pool, and an explicit zero-day candidate
  queue, evaluated with a hold-out zero-day simulation on UNSW-NB15.
\end{enumerate}

\section{Related Work}\label{sec:related}
% TODO: expand from ieee_paper_draft.md Section II once final experiments land.
Isolation Forest \cite{liu2008} and extensions \cite{hariri2021}; LOF
\cite{breunig2000}; autoencoder NIDS \cite{mirsky2018,zavrak2020}; EVT for
anomaly thresholds \cite{siffer2017}; CTU-13 \cite{garcia2014} and UNSW-NB15
\cite{moustafa2015} benchmarks. Our work differs by diagnosing the ceiling
formally before proposing remedies, and by validating each closed-form
prediction on real traffic.

\section{Theoretical Framework}\label{sec:theory}

Let $S$ be a scalar anomaly score (higher $=$ more anomalous),
$Y\in\{0,1\}$ the true class ($1=$ attack), and
$F_A(t)=P(S\le t\mid Y{=}1)$, $F_N(t)=P(S\le t\mid Y{=}0)$ the
class-conditional CDFs. A threshold detector decides ``attack'' iff $S>t$;
its operating point is $\TPR(t)=1-F_A(t)$, $\FPR(t)=1-F_N(t)$, and its Youden
index is $J(t)=\TPR(t)-\FPR(t)$.

\subsection{The KS Ceiling}

\begin{theorem}[KS ceiling]\label{thm:ks}
For any threshold detector on $S$, in either orientation,
\[
\sup_t |J(t)| \;=\; \sup_t |F_N(t)-F_A(t)| \;=\; \KS(F_A,F_N),
\]
and the bound is invariant under any monotone transformation of $S$.
\end{theorem}

\begin{proof}
$J(t) = (1-F_A(t)) - (1-F_N(t)) = F_N(t)-F_A(t)$; take suprema over $t$ and
over the two orientations. Monotone maps preserve both CDFs' level sets, hence
KS.
\end{proof}

\begin{corollary}[Balanced-accuracy ceiling]\label{cor:ba}
$\max_t \mathrm{BA}(t) = \tfrac12\bigl(1+\KS\bigr)$.
\end{corollary}

\begin{corollary}[Recall at a false-alarm budget]\label{cor:budget}
$\TPR(t) \le \FPR(t) + \KS$ for every $t$; hence at budget $\FPR\le\alpha$,
recall is at most $\alpha + \KS$.
\end{corollary}

\begin{corollary}[AUC bound]\label{cor:auc}
$\mathrm{AUC} \le \tfrac12 + \KS$.
\end{corollary}

\emph{Scope.} For arbitrary measurable decision regions the bound is the total
variation distance $\mathrm{TV}\ge\KS$, with equality under monotone
likelihood ratio; deployed detectors are threshold rules, so KS is the
operative ceiling.

\textbf{Empirical validation (CTU-13, held-out test).} The empirical
$\sup_t J(t)$ equals the two-sample KS statistic to machine precision for both
detectors (difference $\le 5.6\times10^{-17}$). Corollary~\ref{cor:ba}:
autoencoder $\KS=0.863$ gives ceiling $0.931$; observed balanced accuracy
$0.929$ --- \emph{the deployed detector operates at its theoretical ceiling}.
Corollary~\ref{cor:budget}: Isolation Forest ($\KS=0.457$) cannot exceed
$50.7\%$ recall at a $5\%$ false-alarm budget, explaining the widely reported
plateau.

\subsection{OR-Fusion of Conditionally Independent Detectors}

Let $B_i=\{S_i>t_i\}$, $i=1,2$, be the flag events of two detectors with
operating points $(\TPR_i,\FPR_i)$, and assume conditional independence (CI)
given each class: $P(B_1\cap B_2\mid Y{=}y)=P(B_1\mid y)P(B_2\mid y)$.

\begin{theorem}[Fusion law]\label{thm:fusion}
Under (CI), the OR-fused detector satisfies:
\begin{enumerate}
  \item[(a)] \emph{Miss-product law:}\quad
  $1-\TPR_{\OR} = (1-\TPR_1)(1-\TPR_2)$;
  \item[(b)] \emph{Subadditive false alarms:}\quad
  $\FPR_{\OR} = 1-(1-\FPR_1)(1-\FPR_2) \le \FPR_1+\FPR_2$;
  \item[(c)] \emph{Fusion-benefit criterion:}\quad
  $J_{\OR} > J_1 \iff (1-\TPR_1)\,\TPR_2 > (1-\FPR_1)\,\FPR_2$.
\end{enumerate}
\end{theorem}

\begin{proof}
(a),(b): apply (CI) to the complements --- a union miss is a simultaneous miss.
(c): $J_{\OR}-J_1 = (1-\TPR_1)\TPR_2-(1-\FPR_1)\FPR_2$ by direct expansion.
\end{proof}

\begin{proposition}[Mixture limit]\label{prop:mixture}
If the attack population is a mixture in which subfamily $k$ is invisible to
the other head (likelihood ratio $\equiv 1$), the Neyman--Pearson optimal
region degenerates to the union $B_1\cup B_2$: OR-fusion is exactly optimal
for heterogeneous attack families.
\end{proposition}

We diagnose (CI) with the miss-redundancy coefficient
$\rho_{\mathrm{miss}} = P(\bar B_1\cap\bar B_2\mid Y{=}1)\,/\,
\bigl[P(\bar B_1\mid Y{=}1)P(\bar B_2\mid Y{=}1)\bigr]$
($\rho=1$: independent blind spots) and the conditional mutual information
$I(S_1;S_2\mid Y)$ estimated by within-class quantile binning.

\textbf{Empirical validation (UNSW-NB15, flow-AE head vs.\ temporal-AE head).}
Head operating points $\TPR_1=0.086$, $\TPR_2=0.440$ (both at
95th-percentile thresholds). The law predicts union recall
$1-(1-0.086)(1-0.440)=0.488$; measured: $0.491$. Overall
$\rho_{\mathrm{miss}}=0.994$ and $I(S_1;S_2\mid Y)=0.28$ bits --- the two
anomaly notions are empirically conditionally independent. Per-family
$\rho_{\mathrm{miss}}\in[1.00,1.13]$. The fusion criterion holds with a
$6\times$ margin. Fused Stage-1 recall on hidden families: Backdoor
$8\%\!\to\!41\%$, Analysis $3\%\!\to\!35\%$.

\subsection{Calibrated Thresholds via Extreme Value Theory}

\begin{theorem}[Pickands--Balkema--de~Haan, applied]\label{thm:evt}
If the normal-score distribution lies in the maximum domain of attraction of
an extreme value law, the excess distribution over a high threshold $u$
converges uniformly to a generalized Pareto distribution
$G_{\xi,\sigma}(x) = 1-(1+\xi x/\sigma)^{-1/\xi}$. Consequently, with
$p_u = P(S_N > u)$ estimated empirically and $(\hat\xi,\hat\sigma)$ fitted by
maximum likelihood on the excesses, the threshold achieving a target
false-positive rate $q$ is
\[
t_q \;=\; u + \frac{\hat\sigma}{\hat\xi}
\Bigl[\bigl(p_u/q\bigr)^{\hat\xi} - 1\Bigr]
\qquad (\hat\xi\neq 0).
\]
\end{theorem}

\textbf{Empirical validation (UNSW-NB15 flow-AE, $n=56{,}000$ training
normals, $u=$ 90th percentile).} The fitted shape is $\hat\xi = 1.12$: the
reconstruction-error tail is so heavy its variance is infinite, which is
precisely why percentile thresholds mis-calibrate. At target $q=10^{-3}$ the
GPD threshold realises $\FPR=4.9\times10^{-4}$ on held-out normals while the
percentile threshold realises $3.8\times10^{-3}$ ($3.8\times$ over budget).
The GPD extrapolates analytically to budgets below the $1/n$ resolution of the
data, removing the last free parameter from Stage~1.

\section{Datasets and Experimental Protocol}\label{sec:data}
% TODO: condense from ieee_paper_draft.md Sec. III (CTU-13, UNSW-NB15, preprocessing, no-leakage protocol).
CTU-13 \cite{garcia2014}: 38{,}898 botnet flows, 53{,}314 normal, 57
CICFlowMeter features; balanced 6k/6k samples, 70/30 stratified split.
UNSW-NB15 \cite{moustafa2015}: 175{,}341 train / 82{,}332 test flows, nine
attack families, 39 Argus features. All preprocessing (median imputation,
$\log(1{+}x)$ on skewed columns, standardisation) fitted on training partitions
only.

\section{The Two-Stage Architecture}\label{sec:system}
Stage~1: dual autoencoder heads (flow features; temporal/behavioural features),
each trained on normal traffic only, fused by OR per
Theorem~\ref{thm:fusion}. Stage~2: gradient-boosted classifier over the
flagged pool, trained on the five \emph{known} attack families only; flows
below confidence $\tau=0.55$ enter the \emph{zero-day candidate queue}.
Four families (Backdoor, Analysis, Shellcode, Worms) are withheld from
Stage~2 at all times.

\subsection{Stage-1 Ablation (UNSW-NB15 test)}
\begin{table}[!t]
\caption{Stage-1 configurations. Union is the proposed head; intersection is
the control.}
\label{tab:ablation}
\centering
\begin{tabular}{lcccc}
\toprule
Config & TPR & FPR & Backdoor & Analysis \\
\midrule
Flow head only        & 0.086 & 0.086 & 0.08 & 0.03 \\
Temporal head only    & 0.440 & 0.073 & 0.39 & 0.35 \\
\textbf{Union (ours)} & \textbf{0.491} & 0.143 & \textbf{0.41} & \textbf{0.35} \\
Intersection (control)& 0.034 & 0.016 & 0.06 & 0.02 \\
\bottomrule
\end{tabular}
\end{table}

Table~\ref{tab:ablation} confirms Theorem~\ref{thm:fusion}: the union recall
(0.491) matches the miss-product prediction (0.488), and the intersection
control collapses as expected.

\subsection{End-to-End Zero-Day Simulation}
\begin{table}[!t]
\caption{Full pipeline, identical Stage~2, only Stage~1 changes.}
\label{tab:pipeline}
\centering
\begin{tabular}{lcc}
\toprule
Metric & Flow-only & Dual-head \\
\midrule
Pipeline recall     & 0.084 & \textbf{0.488} \\
Pipeline precision  & 0.850 & \textbf{0.939} \\
Pipeline $F_1$      & 0.152 & \textbf{0.643} \\
Zero-day queue precision & 0.753 & \textbf{0.867} \\
Backdoor end-to-end routing  & 5.0\% & \textbf{27.8\%} \\
Analysis end-to-end routing  & 2.2\% & \textbf{23.8\%} \\
Shellcode end-to-end routing & 0.5\% & 2.1\% \\
\bottomrule
\end{tabular}
\end{table}

The dual head raises pipeline $F_1$ by $4.2\times$ while \emph{increasing}
precision (Table~\ref{tab:pipeline}). Of the 518 hidden-family flows flagged
by Stage~1, 332 route to the zero-day queue, 182 are alerted as
family-misattributed known attacks, and only 4 ($<1\%$) are cleared as benign:
hidden attacks are misattributed, not missed. Family misattribution is an
honest limitation --- the alert fires, but the analyst receives the wrong
attack label. Shellcode remains largely undetected at Stage~1 (2.1\%): a
single-shot exploit leaves neither a flow-statistical nor a temporal
footprint, delimiting what flow-level detection can achieve.

On CTU-13 the single-head pipeline attains $F_1=0.945$ with an 87\%
false-positive reduction over Stage~1 alone.

\section{Discussion and Limitations}\label{sec:discussion}
% TODO: representation bottleneck (Shellcode ~1% at flow level), dataset vintage, Stage-2 label requirement.

\section{Conclusion}\label{sec:conclusion}
% TODO: write last.

\bibliographystyle{IEEEtran}
\begin{thebibliography}{9}

\bibitem{liu2008}
F.~T. Liu, K.~M. Ting, and Z.-H. Zhou, ``Isolation forest,'' in \emph{Proc.
IEEE ICDM}, 2008, pp. 413--422.

\bibitem{breunig2000}
M.~M. Breunig, H.-P. Kriegel, R.~T. Ng, and J.~Sander, ``{LOF}: Identifying
density-based local outliers,'' in \emph{Proc. ACM SIGMOD}, 2000, pp. 93--104.

\bibitem{garcia2014}
S.~Garcia, M.~Grill, J.~Stiborek, and A.~Zunino, ``An empirical comparison of
botnet detection methods,'' \emph{Computers \& Security}, vol.~45, pp.
100--123, 2014.

\bibitem{moustafa2015}
N.~Moustafa and J.~Slay, ``{UNSW-NB15}: A comprehensive data set for network
intrusion detection systems,'' in \emph{Proc. MilCIS}, 2015, pp. 1--6.

\bibitem{moustafa2016}
N.~Moustafa and J.~Slay, ``The evaluation of network anomaly detection
systems,'' \emph{Information Security Journal}, vol.~25, no. 1--3, pp. 18--31,
2016.

\bibitem{mirsky2018}
Y.~Mirsky, T.~Doitshman, Y.~Elovici, and A.~Shabtai, ``Kitsune: An ensemble of
autoencoders for online network intrusion detection,'' in \emph{Proc. NDSS},
2018.

\bibitem{zavrak2020}
S.~Zavrak and M.~Iskefiyeli, ``Anomaly-based intrusion detection from network
flow features using variational autoencoder,'' \emph{IEEE Access}, vol.~8, pp.
108346--108358, 2020.

\bibitem{hariri2021}
S.~Hariri, M.~C. Kind, and R.~J. Brunner, ``Extended isolation forest,''
\emph{IEEE Trans. Knowl. Data Eng.}, vol.~33, no.~4, pp. 1479--1489, 2021.

\bibitem{siffer2017}
A.~Siffer, P.-A. Fouque, A.~Termier, and C.~Largou\"et, ``Anomaly detection in
streams with extreme value theory,'' in \emph{Proc. ACM KDD}, 2017, pp.
1067--1075.

\end{thebibliography}

\end{document}
